(1) 对周期性盒子的密集动量态，\(k\) 空间单位体积可容纳 \(V/(2\pi)^3\) 个轨道态，自旋简并为2。因此
\[
N=2\frac{V}{(2\pi)^3}\frac{4\pi k_F^3}{3}
=\frac{Vk_F^3}{3\pi^2},
\quad k_F=(3\pi^2n)^{1/3}.
\]
故
\[
E_F=\frac{\hbar^2}{2m}(3\pi^2n)^{2/3},
\qquad v_F=\frac{\hbar k_F}{m}.
\]

(2) 零温每个费米球内状态均被占据：
\[
U=2\frac V{(2\pi)^3}\int_{|\boldsymbol k|<k_F}
\frac{\hbar^2k^2}{2m}\,d^3k
=\frac{V\hbar^2k_F^5}{10\pi^2m}.
\]
用 \(N=Vk_F^3/(3\pi^2)\) 消去 \(V,k_F\) 组合，
\[
U=\frac35NE_F.
\]
此内能纯粹来自 Pauli 排斥导致的动能占据，不需要粒子之间存在排斥势。

(3) 固定 \(N\) 时 \(n=N/V\)、\(E_F\propto V^{-2/3}\)，因此 \(U\propto V^{-2/3}\)，其压力
\[
P=-\left(\frac{\partial U}{\partial V}\right)_N
=\frac23\frac UV=\frac25nE_F.
\]
压力与密度满足 \(P\propto n^{5/3}\)，故
\[
B=-V\left(\frac{\partial P}{\partial V}\right)_N
=n\left(\frac{\partial P}{\partial n}\right)_{S}
=\frac53P=\frac23nE_F>0.
\]
这说明压缩体系需要做正功，即使温度为零。

(4) 零温流体的质量密度为 \(\rho=mn\)；压缩过程中绝热与零温基态一致，故
\[
c_s^2=\left(\frac{\partial P}{\partial\rho}\right)_S
=\frac1m\frac{dP}{dn}
=\frac{2E_F}{3m}=\frac{v_F^2}{3},
\qquad c_s=\frac{v_F}{\sqrt3}.
\]
当 \(E_F\) 不再远小于 \(mc^2\)、\(v_F\) 接近光速，或相互作用/晶格能不可忽略时，这个非相对论自由费米气体模型会失效；不能将公式扩展到任意高密度。
